\noindent\textbf{Elección del procedimiento.} Sin impulso exterior apreciable, el momento total inicial nulo se conserva en el marco de la estación. Elegimos positivo el sentido de la herramienta: $0=m_av_a+m_hv_h$. El signo de la velocidad del astronauta debe resultar negativo.

\begin{enumerate}

\item Pasar el momento de la herramienta al otro lado y dividir por la masa del astronauta.

\[\fitmath{\mathrm{va}=\frac{-\mathrm{mh}\cdot \mathrm{vh}}{\mathrm{ma}}}\]

\[\fitmath{\mathrm{va}=\frac{-(20)\cdot (1.5)}{(80)}\approx -0.375\,\mathrm{m/s}}\]

\item Volver al balance inicial para comprobar que los dos momentos se cancelan.

\[\fitmath{p=\mathrm{ma}\cdot \mathrm{va}+\mathrm{mh}\cdot \mathrm{vh}}\]

\[\fitmath{p=(80)\cdot (-0.375)+(20)\cdot (1.5)\approx 0\,\mathrm{kg\,m/s}}\]

\end{enumerate}

\noindent\textbf{Contraste y comparación.} La celeridad pedida es el módulo $0.375\,\mathrm{m/s}$. No se debe usar la masa conjunta en el denominador: el dato de 1.5 m/s está referido a la estación, no al astronauta.
